\chapter{Justificación}
% \label{chapter:PA}

El incremento acelerado de residuos de aparatos eléctricos y electrónicos está directamente relacionado a la rápida innovación tecnológica, el aumento de la demanda y la ausencia de sistemas efectivos de recuperación y reciclaje. Estudios recientes han documento esta problemática en Bogotá – Colombia, donde se evidencia que la clasificación de estos residuos continúa siendo principalmente manual, destacando la baja tasa de aprovechamiento, la informalidad laboral y la limitada infraestructura de clasificación. En consecuencia, dicha problemática genera retos estructurales y tecnológicos significativos para la separación de residuos reciclables \cite{a.ceronRecyclableWasteClassification2025a}.\\

La ineficiencia en la clasificación de residuos de aparatos eléctricos y electrónicos es el principal “cuello de botella”, pues dentro de un entorno manual evita establecer con precisión el valor económico recuperable, asimismo, evita lograr avances en la economía circular. Por otro lado, la ausencia de sistemas que integren la visión computacional con esquemas de valoración estructurados genera una brecha en la toma de decisiones sobre el destino final de los residuos, un problema que requiere soluciones de inteligencia artificial automatizadas para la identificación precisa \cite{jesiddiazbastosImpactoResiduosAparatos2025}.\\

La necesidad de esta investigación surge de la velocidad a la que aumenta los residuos de aparatos eléctricos y electrónicos siendo esta más veloz que nuestra capacidad de crear sistemas para reciclarlos, convirtiéndose así en un problema de vulnerabilidad ambiental. De esta manera, como lo destaca \cite{janoClosingLoopSystematic2026} sin una gestión circular apoyada en tecnologías digitales, la pérdida de materiales valiosos seguirá aumentando. En comparación con estudios anteriores, actualmente se busca reducir la toxicidad y fomentar el \textit{upcycling} de residuos para la transición energética mediante el uso de inteligencia artificial \cite{dagwar_upcycling_2026}.\\

En el contexto colombiano, la problemática es crítica debido al impacto de estos residuos en el desarrollo sostenible local \cite{jesiddiazbastosImpactoResiduosAparatos2025}. Este estudio tiene como propósito desarrollar un sistema de apoyo a la toma de decisiones, integrando tecnologías de IA que permitan la recuperación eficiente de materiales críticos, reduciendo la dependencia de importaciones y fortaleciendo la economía circular local \cite{souzaCHALLENGESTRANSITIONCIRCULAR2025}.\\

La viabilidad técnica de este sistema se establece en la capacidad de implementar arquitecturas de hardware de bajo consumo y modelos de aprendizaje profundo de alta precisión. Mientras que \cite{p.sLowPowerEdgeAI2026} demuestran que es posible realizar deducción local en cinco (5) segundos mediante el uso de \textit{Edge AI} y fusión multisensorial (peso y metales), el estudio de valida que el uso de redes neuronales convolucionales optimizadas ha alcanzado precisiones superiores al 98\%. La integración de ambos enfoques en esta investigación, permite que el algoritmo de valoración económica propuesto cuente con datos de entrada altamente fiables, transformando la clasificación de RAEE de un proceso manual ineficiente en una operación automatizada, explicable y alineada con los estándares globales de economía circular. Del mismo modo, la disponibilidad de datos como el dataset de \cite{vogelgesangMaterialFlowComposition2026} logra el entrenamiento de modelos de visión sin requerir infraestructuras computacionales restringidas.\\

El desarrollo del sistema propuesto busca beneficiar directamente a los gestores de residuos al estandarizar la identificación inteligente de componentes en centros de acopio, eliminando la variabilidad del criterio humano \cite{varshaDataDrivenEWasteManagement2025a} y maximizando la logística para empresas especializadas en economía circular \cite{dehabadi_how_2026}. Por otra parte, la ciudadanía se beneficia mediante la creación de sistemas de asistencia inteligente que facilitan la disposición correcta de residuos \cite{madusankaUrbanGreenEWasteDetection2025a}.\\

Metodológicamente, la presente investigación aporta integración sistemática entre la clasificación visual y la valoración económica de residuos de aparatos eléctricos y electrónicos. En primer lugar, permite reemplazar la subjetividad propia de procesos manuales por un enfoque basado en \textit{Computer Vision} y \textit{Transfer Learning}. Esto posibilita estandarizar la clasificación bajo métricas de precisión demostrables \cite{p.sLowPowerEdgeAI2026}. Esta estandarización es respaldada por \cite{janoClosingLoopSystematic2026}, quien señala que la IA elimina la imprecisión operativa en la economía circular, y por \cite{donmez_benchmarking_2025}, quien establece que el uso de métricas objetivas de desempeño está orientada a garantizar la pureza y el rendimiento en reaprovechamiento. Otro aporte metodológico, radica en la adopción de caracterización de tipos de materiales como lo menciona \cite{vogelgesangMaterialFlowComposition2026}, un método que permite deducir el valor sin la necesidad de análisis fisioquímicos costosos, lo cual, representa un beneficio en la capacidad de obtener estimaciones económicas significativas en tiempos de clasificación reducidos.
